% Generated draft: requires mathematical and source review.
\documentclass[11pt]{article}
\usepackage[margin=25mm]{geometry}
\usepackage{fontspec}
\setmainfont{Latin Modern Roman}
\usepackage{amsmath,amssymb,mathtools,unicode-math}
\setmathfont{Latin Modern Math}
\usepackage{enumitem,needspace,xurl,hyperref}
\setlength{\parindent}{0pt}
\setlength{\parskip}{5pt}
\setlength{\emergencystretch}{3em}
\newcommand{\N}{\mathbb{N}}
\newcommand{\Z}{\mathbb{Z}}
\newcommand{\Q}{\mathbb{Q}}
\newcommand{\R}{\mathbb{R}}
\newcommand{\C}{\mathbb{C}}
\newcommand{\F}{\mathbb{F}}
\newcommand{\A}{\mathbb{A}}
\newcommand{\PP}{\mathbb{P}}
\newcommand{\E}{\mathbb{E}}
\newcommand{\eps}{\varepsilon}
\newcommand{\dd}{\,\mathrm{d}}
\begin{document}
\section*{UnsolvedMath ID: 10000027}
\textbf{Source problem code:} catalogue entry
\subsection{Definitions and mathematical statement}
Let \(G_n=(V_n,E_n)\) be finite graphs with \(|V_n|\to\infty\), maximum degree at most \(D\), and uniform edge expansion: for some \(h>0\),
\[
|\partial_E A|\ge h|A|
\quad\text{whenever }0<|A|\le|V_n|/2.
\]
The source assumes local convergence to an infinite graph \(G\), without specifying roots. In a usual local-weak model a root \(o_n\) is uniform in \(V_n\), and the rooted neighborhoods converge in distribution to those of \((G,o)\); an almost-sure fixed rooted limit is a stronger specification than an arbitrary random rooted limit.

For independent bond percolation on a connected infinite graph, a standard critical value is
\[
p_c(G)=\inf\{p\in[0,1]:\mathbb P_p(o\leftrightarrow\infty)>0\}.
\]
The intended assertion is
\[
\lim_{n\to\infty}p_c(G_n)=p_c(G).
\]
But the infinite-cluster definition does not define a finite-graph threshold. A finite giant-component threshold needs extra conventions. For example,
\[
p_{n,\alpha}=\inf\{p:\mathbb P_p(|C_{\max}(G_n)|\ge\alpha|V_n|)\ge1/2\}
\]
depends on \(\alpha>0\), and holding a positive fraction fixed is not automatically the onset threshold corresponding to \(p_c(G)\). The source must specify whether a vanishing fraction, a double limit, another sharp-threshold convention, or a different finite observable defines its \(p_c(G_n)\). Site versus bond percolation is also not stated explicitly; bond percolation is displayed as the usual model.

No transitivity, regularity, or deterministic-limit assumption is silently added. The supplied research summary reports major locality progress but expressly says its latest-status audit was incomplete. Results for vertex-transitive graphs or specially controlled local limits do not by themselves establish the stated arbitrary-expander version. The record provides no precise modern theorem resolving this exact finite-threshold formulation.
\subsection{Short English statement}
Do percolation thresholds converge along bounded-degree expanders with an infinite local limit? Specify the finite giant-component threshold, local convergence/root law, and percolation model before asserting the limit.
\subsection{Sources}
\begin{itemize}
\item Local UnsolvedMath registry, permanent ID \(10000027\), exact source code \texttt{catalogue entry}, ``Percolation thresholds along expander limits.'' Proposer Itai Benjamini.
\item Benjamini, \emph{Coarse Geometry and Randomness}, Saint-Flour notes, Open problem 5.10, recorded PDF page 40; \url{https://www.wisdom.weizmann.ac.il/~itai/stflouraug24.pdf#page=40}. Recorded access 2026-07-29; extracted text checked against rendered source.
\item Recorded preservation copy: \url{https://arquivo.pt/noFrame/replay/20201231041548id_/http://www.wisdom.weizmann.ac.il/~itai/stflouraug24.pdf}.
\item Benjamini--Nachmias--Peres, ``Does the percolation threshold depend on the graph locally?,'' \emph{Annals of Probability} (2011), arXiv:0910.1809, as supplied. The summary also mentions Easo/Hutchcroft and an uncertain ``Ph. Souplet-type'' attribution in 2022--2025, without complete references or exact matching hypotheses.
\end{itemize}
\end{document}
